\subsection{Freely homotopic loops}

DEFINITION 4. --- Let X be a topological space and let c and c' be two loops in X. One calls a free homotopy joining c to c' a homotopy σ joining c to c' such that σ(0, s) = σ(1, s) for all s ∈ I. One says that c is freely homotopic to c' if there exists a free homotopy joining c to c'.

The free homotopies joining \(c\) to \(c'\) correspond to the paths in the space \(\Omega(\mathrm{X})\) of loops of X joining \(c\) to \(c'\) . Consequently, the relation “ \(c\) is freely homotopic to \(c'\) ” is an equivalence relation on \(\Omega(\mathrm{X})\) whose equivalence classes are the arcwise connected components of \(\Omega(\mathrm{X})\) .

Note. --- Let \(\varphi\) be the canonical map from \(\mathbf{R}\) onto \(\mathbf{T} = \mathbf{R}/\mathbf{Z}\) (TG, V, p.~2). The map \(f \mapsto f \circ \varphi|\mathbf{I}\) is a homeomorphism of \(\mathcal{C}_c(\mathbf{T};\mathbf{X})\) onto \(\Omega(\mathbf{X})\) , whence, by passing to arcwise connected components, a bijection from the set {[}T; X{]} (III, p.~230) onto the set of free homotopy classes of loops in X.

PROPOSITION 6. --- Let X be an arcwise connected topological space and let x be a point of X.

\begin{enumerate}
\def\labelenumi{\alph{enumi})}
\item
  Every loop in X is freely homotopic to a loop at x. More precisely, if c is a loop at y and d is a path with origin y and endpoint x, c is freely homotopic to the loop \((\overline{d} * c) * d\) at x.
\item
  Two loops in X at x are freely homotopic if and only if their strict homotopy classes are conjugate in the group \(\pi_{1}(X,x)\) .
\end{enumerate}

Let us prove \(a\) ). For every \(s \in [0,1]\) , denote by \(d_s\) the path in X defined by \(d_s(t) = d(st)\) for \(t \in \mathbf{I}\) ; its origin is \(y\) . As the map \((s,t) \mapsto d(st)\) is continuous, the map \(s \mapsto d_s\) from \(\mathbf{I}\) to \(\mathcal{C}_c(\mathbf{I};\mathbf{X})\) is continuous (III, p.~257, prop. 1). The map \(s \mapsto (\overline{d_s} * c) * d_s\) is then a path in \(\Omega(\mathbf{X})\) (III, p.~257, prop. 2) joining \((e_y * c) * e_y\) to \((\overline{d} * c) * d\) , whence \(a\) ).

Let c and \(c'\) be two loops in X at x. If their strict homotopy classes are conjugate in \(\pi_{1}(X,x)\) , there exists a loop d at x such that \(c'\) is strictly homotopic to the loop \((\overline{d} * c) * d\) . It follows from a) that c and \(c'\) are freely homotopic. Conversely, suppose there exists a free homotopy \(\varphi\) joining c to \(c'\) . Set \(d(t) = \varphi(0,t)\) , we also have \(d(t) = \varphi(1,t)\) and d is a loop at x. By lemma 1 of III, p.~295, the loops \(c * d\) and \(d * c'\) are strictly homotopic. The strict homotopy classes of c and \(c'\) are therefore conjugate in \(\pi_{1}(X,x)\) .

Scholium. --- Let X be an arcwise connected topological space and let x be a point of X. Proposition 6 allows one to define a canonical bijection from the set of free homotopy classes of loops in X onto the set of conjugacy classes in \(\pi_{1}(X,x)\) .

